\documentclass[11pt]{article}
\usepackage[margin=1in]{geometry}
\usepackage{amsmath,amssymb,amsthm,mathtools}
\usepackage{graphicx}
\usepackage{xcolor}
\usepackage{microtype}
\usepackage[hidelinks]{hyperref}

\newtheorem{theorem}{Theorem}
\newtheorem{proposition}[theorem]{Proposition}
\newtheorem{lemma}[theorem]{Lemma}
\newcommand{\D}{\mathbb D}
\newcommand{\N}{\mathbb N}
\newcommand{\St}{\operatorname{Stolz}}
\newcommand{\Hinfty}{H^\infty}
\newcommand{\Cay}{\mathcal C}
\newcommand{\SM}{\mathcal S}

\title{Coefficient spaces and an explicit Riemann map\\
for the Stolz-domain question in arXiv:2303.03022}
\author{Autonomous mathematical investigation, lane 5}
\date{13 August 2026}

\begin{document}
\maketitle

\begin{abstract}
The open-question subsection of Fackler and Hauer asks for the spaces generated
by products of their standard Ritt square functions, for a general
fixed-coefficient representation, and for an explicit conformal map from
their Stolz domain to the disc.  We identify every product space exactly as a
weighted analytic coefficient space and identify their algebraic sum.  This
also gives a sharp obstruction: the proposed products cannot represent every
bounded analytic function, because all their finite sums extend to the whole
unit disc.  We then give and prove an explicit normalized Riemann map in terms
of a generalized Chebyshev function.  Finally, we show why simply pulling
standard monomial square functions back through this map cannot solve the
remaining square-function problem.  Thus the coefficient-space and conformal
map parts are complete; the request for a genuinely new elementary square
family remains open.
\end{abstract}

\section{Question and answer}

For $\omega>1$, set
\[
  \St_\omega=\left\{z\in\D:
  \frac{|1-z|}{1-|z|}<\omega\right\}.
\]
The source defines, for integers $m\geq1$ and $k\geq1$,
\[
  \varphi_m(k,z)=k^{m-\frac12}(1-z)^m z^{k-1}
\]
and asks to identify
\[
 \SM_{m_1,m_2}=
 \left\{\sum_{k\geq1}a_k\varphi_{m_1}(k,z)
 \varphi_{m_2}(k,z):(a_k)\in\ell_\infty\right\},
\]
their sum in $\Hinfty(\St_\omega)$, and the conformal map of
$\St_\omega$ onto $\D$.

For an integer $s\geq0$, introduce the weighted coefficient space
\[
 \mathcal A_s=\left\{A(w)=\sum_{n\geq0}c_nw^n:
      \sup_{n\geq0}\frac{|c_n|}{(n+1)^s}<\infty\right\}.
\]
Also write $D=w\frac{d}{dw}$ and
\[
 \Hinfty_{\rm coeff}=\left\{\sum_{n\geq0}b_nw^n:
 (b_n)\in\ell_\infty\right\}.
\]

\begin{theorem}[Complete coefficient and conformal-map subresults]
\label{thm:main}
Let $m_1,m_2\geq1$ and $M=m_1+m_2$.
\begin{enumerate}
\item The space depends only on $M$, and exactly
\begin{align*}
 \SM_{m_1,m_2}=\SM_M
 &=\left\{(1-z)^M A(z^2):A\in\mathcal A_{M-1}\right\} \\
 &= (1-z)^M(D+1)^{M-1}\Hinfty_{\rm coeff}(z^2).
\end{align*}
Consequently the algebraic sum requested in the source is
\[
 \sum_{m_1,m_2\geq1}^{\rm alg}\SM_{m_1,m_2}
 =\left\{\sum_{M\in F}(1-z)^M A_M(z^2):
 F\subset\{2,3,\ldots\}\text{ finite},\ A_M\in\mathcal A_{M-1}\right\}.
\]
\item This algebraic sum is a proper subspace of
$\Hinfty(\St_\omega)$.  In particular, the fixed-coefficient formula
proposed as (6.5) in the source cannot represent all of
$\Hinfty(\St_\omega)$, for any choices of $m_1,m_2$.
\item Put
\[
 a=\frac1\omega,\qquad \beta=\arccos a,\qquad
 \kappa=\frac{\pi}{2\beta},\qquad C(z)=\frac{1+z}{1-z}.
\]
Let the generalized Chebyshev function be the branch
\[
 T_\kappa(w)={}_2F_1\left(-\kappa,\kappa;\frac12;
                     \frac{1-w}{2}\right),
\]
analytic on $\mathbb C\setminus(-\infty,-1]$.  Equivalently,
$T_\kappa(\cosh u)=\cosh(\kappa u)$ whenever
$|\operatorname{Im}u|<\beta$.  Then
\begin{equation}
 \boxed{\quad
 F_\omega(z)=
 \frac{T_\kappa(C(z))-1}{T_\kappa(C(z))+1}
 \quad}
 \label{eq:map}
\end{equation}
is a conformal bijection from $\St_\omega$ onto $\D$.  It is normalized by
\[
 F_\omega(0)=0,\qquad F_\omega'(0)=\kappa^2>0,\qquad
 F_\omega\!\left(\frac{1-\omega}{1+\omega}\right)=-1,\qquad
 \lim_{z\to1}F_\omega(z)=1.
\]
\end{enumerate}
\end{theorem}

\section{Idea of the proof}

The coefficient question collapses immediately after multiplying the two
standard functions: the result is an even power series, multiplied by
$(1-z)^M$, whose $n$th coefficient has exactly the allowed growth
$(n+1)^{M-1}$.  This not only identifies the spaces but exposes their common
analytic-continuation property.  A single pole placed in the part of the unit
disc outside the closed Stolz domain then rules out the hoped-for universal
fixed-coefficient representation.

For the conformal map, the Cayley transform turns the defining metric ratio
into a difference of distances from the foci $\pm1$.  Hence the boundary
becomes a hyperbola, parametrized by $\cosh(t\pm i\beta)$.  The hyperbola
domain is the quotient of the strip $|\operatorname{Im}u|<\beta$ by
$u\sim-u$.  Scaling this strip to $|\operatorname{Im}v|<\pi/2$ and applying
$\cosh$ again maps the quotient to the right half-plane.  This is precisely a
generalized Chebyshev function; a final Cayley map gives the disc formula.

\section{The coefficient spaces}

The product defining the source space is
\begin{equation}
 \varphi_{m_1}(k,z)\varphi_{m_2}(k,z)
 =k^{M-1}(1-z)^M z^{2k-2}.
 \label{eq:product}
\end{equation}
Thus, if $n=k-1$, the coefficient of $w^n$ in the corresponding $A(w)$
is $(n+1)^{M-1}a_{n+1}$.  This proves both set equalities in part 1,
because
\[
 (D+1)^{M-1}\sum_{n\geq0}b_nw^n
 =\sum_{n\geq0}(n+1)^{M-1}b_nw^n.
\]
Conversely, every $A\in\mathcal A_{M-1}$ produces a bounded function on
$\St_\omega$.  Indeed, for $r=|z|$,
\begin{align*}
 |(1-z)^MA(z^2)|
 &\leq C\|A\|_{\mathcal A_{M-1}}|1-z|^M
        \sum_{n\geq0}(n+1)^{M-1}r^{2n}\\
 &\leq C_M\|A\|_{\mathcal A_{M-1}}
       \left(\frac{|1-z|}{1-r^2}\right)^M
 \leq C_M\omega^M\|A\|_{\mathcal A_{M-1}}.
\end{align*}
Here $|1-z|<\omega(1-r)$ and $1-r^2\geq1-r$.
Taking finite sums proves the asserted formula for the algebraic sum.

There is also a decisive obstruction to a universal representation.  Every
$A\in\mathcal A_s$ has radius of convergence at least one, hence every
member of the algebraic sum extends holomorphically to all of $\D$.  On the
other hand, set
\[
 z_0=-\frac{\omega}{\omega+1}.
\]
This point lies in $\D\setminus\overline{\St_\omega}$, so
$f(z)=(z-z_0)^{-1}$ belongs to $\Hinfty(\St_\omega)$.  If $f$ agreed on
$\St_\omega$ with a function $G$ holomorphic on $\D$, the identity theorem
on the connected punctured disc $\D\setminus\{z_0\}$ would give
$G(z)=(z-z_0)^{-1}$ there.  Continuity at $z_0$ would then force
$0=(z-z_0)G(z)|_{z=z_0}=1$, a contradiction.  This proves part 2.

Notice that the conclusion concerns the algebraic sum asked for in the
source; no assertion about its norm closure is needed.

\section{Proof of the explicit conformal map}

The Cayley map $C$ sends $\D$ to the right half-plane.  If $w=C(z)$, then
\[
 z=\frac{w-1}{w+1},\qquad
 \frac{|1-z|}{1-|z|}=\frac{2}{|w+1|-|w-1|}.
\]
Therefore $C(\St_\omega)$ is the component to the right of the hyperbola
\[
 \frac{x^2}{a^2}-\frac{y^2}{1-a^2}=1,\qquad a=\frac1\omega.
\]
Writing $\beta=\arccos a$, its two boundary arcs have the parametrization
\[
 w=\cosh(t\pm i\beta)
   =a\cosh t\ \pm i\sqrt{1-a^2}\sinh t,\qquad t\in\mathbb R.
\]

Let $S_\beta=\{u:|\operatorname{Im}u|<\beta\}$.  The map
$u\mapsto\cosh u$ maps $S_\beta$ onto $C(\St_\omega)$; its only
identification inside this strip is $u\sim-u$.  Similarly,
$v\mapsto\cosh v$ maps $S_{\pi/2}$ onto the right half-plane with the same
identification.  Multiplication by
$\kappa=\pi/(2\beta)$ is a conformal bijection
$S_\beta\to S_{\pi/2}$ and respects $u\sim-u$.  It therefore descends to a
conformal bijection
\[
 w=\cosh u\longmapsto \cosh(\kappa u)=T_\kappa(w)
\]
from $C(\St_\omega)$ onto the right half-plane.

This quotient argument also resolves the apparent branch at $w=1$: locally
$w-1=u^2/2+O(u^4)$ and
$T_\kappa(w)-1=\kappa^2u^2/2+O(u^4)$, so
$T_\kappa'(1)=\kappa^2$.  The hypergeometric expression in the theorem is
the single-valued analytic continuation across $(-1,1)$ and has its first
branch obstruction only at $(-\infty,-1]$, well outside the hyperbola
domain.  Finally, $(q-1)/(q+1)$ maps the right half-plane onto $\D$.
Composing these three bijections gives \eqref{eq:map} and the stated
normalizations.

An explicit inverse, with the strip branch of $\operatorname{arccosh}$, is
\[
 F_\omega^{-1}(\zeta)=
 \frac{\cosh\!\left(\kappa^{-1}
       \operatorname{arccosh}\frac{1+\zeta}{1-\zeta}\right)-1}
      {\cosh\!\left(\kappa^{-1}
       \operatorname{arccosh}\frac{1+\zeta}{1-\zeta}\right)+1}.
\]
Changing the permitted inverse branch only changes the argument by a sign,
which disappears under $\cosh$.

\section{Why the map does not by itself solve the square-function problem}

One tempting upgrade is to replace $z$ in the standard family by
$\zeta=F_\omega(z)$.  It fails at the level of the required pointwise
$\ell_2$ bound.  For $m\geq1$,
\[
 \sum_{k\geq1}k^{2m-1}|1-\zeta|^{2m}|\zeta|^{2k-2}
 \asymp_m
 \left(\frac{|1-\zeta|}{1-|\zeta|^2}\right)^{2m}.
\]
As $\zeta$ approaches any point of $\partial\D\setminus\{1\}$, the right
side diverges.  The conformal map sends the whole curved Stolz boundary to
$\partial\D$, so the cancellation that made the original Ritt family
bounded only at its distinguished vertex is lost.

There is a more general diagonal obstruction.  Suppose two disc families
have the form
\[
 \psi_j(\zeta)=u_j\zeta^{p_j},\qquad
 \widetilde\psi_j(\zeta)=v_j\zeta^{q_j},
 \qquad p_j,q_j\geq0,
\]
and both are uniformly pointwise square summable.  Letting $r\uparrow1$
and using monotone convergence gives $u,v\in\ell_2$, hence
$(u_jv_j)\in\ell_1$.  Every bounded-coefficient sum
$\sum_j a_j\psi_j\widetilde\psi_j$ therefore has absolutely summable Taylor
coefficients, even when exponents repeat.  It lies in the Wiener algebra and
extends continuously to $\overline\D$.  This cannot cover $\Hinfty(\D)$:
for example
\[
 S(\zeta)=\exp\!\left(-\frac{1+\zeta}{1-\zeta}\right)
\]
is bounded and analytic on $\D$ but has no continuous extension at $1$.
Thus a successful explicit family must be genuinely non-diagonal and
localized simultaneously in scale and boundary position.

There is a constructive, but not elementary, fallback.  The later paper
arXiv:2410.22006 constructs contour atoms $\Phi_j$ with
$\sup_z\sum_j|\Phi_j(z)|^{1/2}<\infty$ and bounded-coefficient
Franks--McIntosh representations.  Formula \eqref{eq:map} makes their
canonical factors specifiable: factor
$\Phi_j\circ F_\omega^{-1}=I_jO_j$ in the disc and set
\[
 \psi_j=(I_jO_j^{1/2})\circ F_\omega,\qquad
 \widetilde\psi_j=O_j^{1/2}\circ F_\omega.
\]
The outer factor has the standard boundary-integral formula, and the maximum
principle transfers the $\sum|\Phi_j|^{1/2}$ bound to both factor families.
This supplies explicit coordinates for the known contour construction, but
does not produce the simple Ritt-adapted square family sought in the source.

\section{Verification and literature boundary}

The accompanying script independently checks \eqref{eq:product}, boundary
moduli, the strip formula, the inverse, interior containment, and the
removable derivative at $w=1$ for $\omega=1.2,2,5$.  Boundary and inverse
errors were below $3\times10^{-15}$; the finite-difference derivative errors
at step $10^{-6}$ were below $8\times10^{-6}$.

A targeted search through 13 August 2026 found the 2012/2014 standard Ritt
square-function theory and the 2024 Ritt$_E$ contour decomposition, but no
source giving the exact generalized-Chebyshev map \eqref{eq:map} or the
coefficient-space obstruction above.  This is evidence about search coverage,
not a claim of bibliographic priority.

\begin{thebibliography}{9}
\bibitem{FH}
S. Fackler and M. Hauer,
\emph{Remarks on Ritt operators, their $H^\infty$ functional calculus and
associated square function estimates}, arXiv:2303.03022 (2023).

\bibitem{LM}
C. Le Merdy,
\emph{$H^\infty$ functional calculus and square function estimates for Ritt
operators}, Rev. Mat. Iberoam. 30 (2014), 1149--1190;
arXiv:1202.0768.

\bibitem{B}
O. Bouabdillah,
\emph{Square functions associated with Ritt$_E$ operators},
arXiv:2410.22006 (2024).
\end{thebibliography}

\begin{figure}[htbp]
\centering
\includegraphics[width=0.76\textwidth]{figures/open_question_crop.png}
\caption{The exact open-question subsection on PDF page 15 of
arXiv:2303.03022.}
\end{figure}

\end{document}
